\subsection{正定自旋几何与 Dirac 算子的平方}

Lorentz 场表示是局部结构；全局自旋丛还受拓扑限制。以下使用正定黎曼度规，建立与前面联络拉普拉斯相接的 Dirac 平方恒等式。

对高自旋的数学讨论若只处理整数阶对称张量，仍无法覆盖费米场。在任意空间上引入 Dirac 场不是增加一个矩阵指标那么简单：必须先检查流形是否允许自旋结构，再在自旋丛上构造联络与 Clifford 乘法。曲率由此以精确的 $\mathcal R/4$ 形式进入 Dirac 算子的平方。

设 $M$ 定向且带黎曼度规，其正交标架丛结构群为 $SO(d)$。半整数自旋表示一般属于双覆盖 $\operatorname{Spin}(d)$，而不是 $SO(d)$。要定义全局自旋场，必须存在一个 $\operatorname{Spin}(d)$ 主丛及到正交标架丛的二重覆盖。拓扑障碍为 $w_2(TM)$；若该类非零，则不能选全局自旋结构，尽管局部 Clifford 代数与局部 Dirac 矩阵仍可书写。

若自旋结构存在，取 Clifford 模 $S$，构造自旋丛 $\mathbb S\to M$。局部正交标架 $e_a$ 中的矩阵满足
\begin{equation}\label{eq:ghs-riemannian-clifford}
 \{\gamma^a,\gamma^b\}=2\delta^{ab}\id,
\end{equation}
并由自旋联络 $\nabla_a^{\mathbb S}$ 满足 $\nabla^{\mathbb S}\gamma=0$。

取 Hermitian Clifford 矩阵，定义形式自伴的欧氏 Dirac 算子
\begin{equation}\label{eq:ghs-geometric-dirac}
 \mathsf D=-\ii\gamma^a\nabla_a^{\mathbb S}.
\end{equation}
对于无边界紧致自旋流形，在标准 Sobolev 定义域上该算子自伴；有边界时需另选椭圆边界条件。

\begin{theorem}[Lichnerowicz 公式]
\label{thm:ghs-lichnerowicz-spinor}
在上述 Clifford 约定下，
\begin{equation}\label{eq:ghs-dirac-square-curvature}
 \mathsf D^2=(\nabla^{\mathbb S})^*\nabla^{\mathbb S}+\frac14\mathcal R_g.
\end{equation}
\end{theorem}
\begin{proof}
选一点处的正规正交标架，使该点 $\nabla_{e_a}e_b=0$。由 Clifford 关系，把算子平方分成指标对称与反对称部分：
\[
 \mathsf D^2=-\delta^{ab}\nabla_a^{\mathbb S}\nabla_b^{\mathbb S}
 -\frac14[\gamma^a,\gamma^b][\nabla_a^{\mathbb S},\nabla_b^{\mathbb S}].
\]
四指标仍取 $R_{abcd}=g(R(e_a,e_b)e_c,e_d)$，与前面几何章节一致。局部自旋联络写为 $\nabla^{\mathbb S}=\dd+\tfrac14\omega_{cd}\gamma^c\gamma^d$，其中 $\omega_{cd}=g_{ce}\omega^e{}_d$。其曲率来自同一个 Cartan 第二结构方程，因此
\begin{align*}
 [\nabla_a^{\mathbb S},\nabla_b^{\mathbb S}]
 &=\frac14g_{ce}R^e{}_{dab}\gamma^c\gamma^d\\
 &=\frac14R_{abdc}\gamma^c\gamma^d
 =-\frac14R_{abcd}\gamma^c\gamma^d.
\end{align*}
最后一步使用第三、第四指标反对称性。代入 Dirac 平方，曲率贡献因而是 $+\tfrac18R_{abcd}\gamma^a\gamma^b\gamma^c\gamma^d$。
令 $\gamma^{ab}=\gamma^{[a}\gamma^{b]}$，$\gamma^{abcd}=\gamma^{[a}\gamma^b\gamma^c\gamma^{d]}$，用 Clifford 关系逐项拆出
\begin{align*}
\gamma^a\gamma^b\gamma^c\gamma^d
={}&\gamma^{abcd}
+\delta^{ab}\gamma^{cd}-\delta^{ac}\gamma^{bd}+\delta^{ad}\gamma^{bc}\\
&+\delta^{bc}\gamma^{ad}-\delta^{bd}\gamma^{ac}+\delta^{cd}\gamma^{ab}\\
&+\delta^{ab}\delta^{cd}-\delta^{ac}\delta^{bd}+\delta^{ad}\delta^{bc}.
\end{align*}
第一 Bianchi 恒等式令 $R_{abcd}\gamma^{abcd}=0$。二重 Clifford 部分收缩为对称 Ricci 张量与反对称 $\gamma^{ab}$ 的配对，故为零。在相同曲率号约定下，最后三个标量项给出
\[
 R_{abcd}(-\delta^{ac}\delta^{bd}+\delta^{ad}\delta^{bc})=2\mathcal R_g.
\]
这两个缩并分别是 $-\sum_{a,b}R_{abab}$ 与 $\sum_{a,b}R_{abba}$；按 $\mathcal R_g=\sum_{a,b}R_{abba}$ 及反对称性，各给出一次 $\mathcal R_g$。因此曲率贡献是 $(1/8)(2\mathcal R_g)=\mathcal R_g/4$。对任意点成立，再把正规标架中的 $-\delta^{ab}\nabla_a\nabla_b$ 识别为联络 拉普拉斯算子 $\nabla^*\nabla$，得到 \cref{eq:ghs-dirac-square-curvature}。
\end{proof}

由 \cref{eq:ghs-dirac-square-curvature} 与无边界积分分部，
\[
 \norm{\mathsf D\psi}^2=\norm{\nabla^{\mathbb S}\psi}^2
 +\frac14\int_M\mathcal R_g|\psi|^2\dd\mu_g.
\]
若 $\mathcal R_g>0$ 处处成立，则 $\mathsf D\psi=0$ 只能给 $\psi=0$。这是一个曲率直接约束费米零模的完整几何结论。注意这里 $\mathcal R>0$ 指所选完整流形上的正标量曲率，而不是某个局部截面的单一常数。

整数自旋张量满足对称、无迹条件；张量—自旋量则可考虑 $\operatorname{Sym}^S(T^*M)\otimes\mathbb S$ 的截面，并进一步施加 $\gamma$-迹约束来选定不可约 $\operatorname{Spin}(d)$ 表示。它们的协变微分方程可涉及 Rarita--Schwinger 型算子，因曲率与规范约束而受到更严格的背景限制。

局部 Clifford 表示、几何 Dirac 算子和相互作用高自旋费米场是三种不同层次。本章建立的是前两者与高自旋表示的数学入口；特定的质量、规范冗余与相互作用一致性必须另有动力学输入，不能仅由 \cref{eq:ghs-riemannian-clifford} 猜出。
